\section{Sonuç ve Gelecek Çalışmalar}
Bu raporda, fundus retina görüntülerinde damar segmentasyonu ve segmentasyon çıktısından morfolojik damar göstergelerinin çıkarımı ele alınmıştır. FOV maskeleme tabanlı ön işleme ile retina dışı alanların etkisi azaltılmış; U-Net tabanlı model BCE ve Dice kaybı kombinasyonu ile eğitilmiştir. Test kümesinde macro ortalama Dice $0.8059\pm0.0752$ ve IoU $0.6807\pm0.0946$ elde edilmiştir. Segmentasyon sonrası analiz, damar yoğunluğunun ortalamada korunabildiğini; ancak iskelet uzunluğu ve dallanma sayısı gibi topolojik göstergelerde azalma eğilimi olabileceğini göstermiştir.

Gelecek çalışmalar kapsamında aşağıdaki iyileştirmeler planlanabilir:
\begin{itemize}
    \item Daha güçlü mimariler (ör. attention tabanlı U-Net varyantları) ile ince damarların daha iyi yakalanması,
    \item Kayıp fonksiyonu ve eşikleme stratejilerinin (ör. adaptif eşik, farklı loss kombinasyonları) sistematik biçimde araştırılması,
    \item Segmentasyon sonrası post-process adımları ile topolojik bütünlüğün artırılması ve morfolojik ölçümlerin daha kararlı hâle getirilmesi.
\end{itemize}
